% =====================================================================
%  Harald Høffding, On Pessimism
%  Translated from the Danish: „Om Pessimismen“, Vor Ungdom. Tidsskrift for Opdragelse og Undervisning, 1879, pp. 69--90.
%
%  TRANSLATION -- GENERATED; DO NOT EDIT BY HAND.  Built by ebook-method/hoeffding-1879-82/tr_parts.py build om-pessimismen-1879 --write
%  from the parts in ebook-method/hoeffding-1879-82/tr/ (GLOSSARY.md, TRANSLATOR-BRIEF.md).  Made from the collated
%  transcription texts/hoeffding/om-pessimismen-1879/transcription.tex (2026-10-15); [N] in the margin = the page of the original.
%  Terms: GLOSSARY.md (choices marked ◆ there for Hans to review).
% =====================================================================
\documentclass[12pt,a4paper,openany]{book}
\usepackage[utf8]{inputenc}
\usepackage[T1]{fontenc}
\usepackage{libertinus}
\usepackage{libertinust1math}
\usepackage[german,french,english]{babel}
\usepackage{textalpha}
\usepackage[protrusion=true,expansion=true]{microtype}
\usepackage{setspace}
\usepackage[margin=1.3in]{geometry}
\usepackage{fancyhdr}
\usepackage{hyperref}
\hypersetup{pdftitle={On Pessimism}, pdfauthor={Harald Høffding}, hidelinks}
\pagestyle{fancy}
\fancyhf{}
\fancyhead[LE,RO]{\thepage}
\fancyhead[C]{\textit{H. Høffding: On Pessimism (1879)}}
\fancypagestyle{plain}{\fancyhf{}\fancyhead[LE,RO]{\thepage}}
\setstretch{1.3}
\usepackage{marginnote}
\usepackage{graphicx}
\DeclareUnicodeCharacter{0254}{\reflectbox{c}}   % ɔ (id est)
\renewcommand{\thefootnote}{\textsuperscript{\arabic{footnote}})}
\newcommand{\opage}[1]{\marginnote{\footnotesize\textit{[#1]}}}
\newcommand{\signature}[1]{\par\smallskip\hfill #1\hspace*{2em}\par}
\pdfstringdefDisableCommands{\renewcommand{\footnote}[2][]{}\renewcommand{\opage}[1]{}}

\begin{document}

\begin{titlepage}
\begin{center}
{\Huge \emph{Vor Ungdom}}\\[0.4em]{\large [Our Youth]}\\[0.6em]{\large A Journal of Upbringing and Instruction}\\[0.6em]{\small edited by Jean Pio, H. Trier and P. Voss}\\[0.6em]1879\\[0.4em]{\small Copenhagen: P. G. Philipsen}\\[2em]
{\small Translated from the Danish}
\end{center}
\end{titlepage}
\clearpage

% --- p. 69 ---
\setcounter{footnote}{0}%
\opage{69}\begin{center}{\Large\sffamily\bfseries \emph{On Pessimism.}}\end{center}\medskip

It is with mental health as with bodily health: it shows itself in the fact that one gets no occasion at all to think of it. Health is merely the necessary foundation for full devotion to, and absorption in, activity and work. And since happiness presumably consists essentially in such a full and unhindered living and working, it is always for the most part unconscious. The vital force unfolds itself without one's having any reason to stop and reflect on one's state or one's feeling. Even if there is resistance and there are obstacles, --- so long as they stand in no disproportion to the force that is at hand, they will precisely contribute to increasing and heightening the unfolding of life. Most people probably lead their lives under such conditions. There is in their nature no such deep rift, no such gnawing doubt, as could hinder the instinct of life and the drive of self-preservation from continuing their activity. Their working, enjoying and suffering is determined by a small, surveyable circle of circumstances and tasks, and they have no reason to pass any judgment on life as a whole. Even if single great and deep sorrows in some cases cast dark shadows over larger or smaller stretches of life,
% --- p. 70 ---
\setcounter{footnote}{0}%
\opage{70}they are nevertheless unable to damage the feeling of life permanently in its innermost root.

In various ways the optimistic tendency lying in the instinct of life itself and the drive of self-preservation may be broken. --- Misfortune and suffering may overwhelm the human power of resistance and make free, harmonious development impossible. The individual human being to whom this happens will feel life as a burden and a pain, and he will perhaps extend his judgment to human life in general. Yet a deep and common feeling of life as unhappy will not arise except on the basis of strong and deeply rooted sympathy. So long as misfortune strikes only the individual, or the individuals each in his own way, it will still stand as an accident. But when a human being takes on the sin and misfortune of others through sympathy and compassion, evil begins to gnaw at the root of life. Hence we also see that wherever the conception of life as an evil spreads, a mighty upswing of the sympathetic feelings has taken place at the same time. Buddhism is not only the religion of pessimism but also the religion of a deep, heartfelt and all-embracing compassion. Christianity, according to which the world lies in evil, is at the same time the religion of universal love of humanity. In modern pessimistic philosophy we find --- in theory at any rate --- the two moments likewise united. Schopenhauer taught that the world was the worst of all possible worlds, and at the same time he set up compassion as the principle of morality.

A pessimistic conception of life therefore always presupposes a comparatively high mental development. Without memory, without wide experience, without the ability to put oneself in the place of other feeling beings, the painful feeling cannot reach any high degree of strength and cannot come to determine the value that is ascribed to life. The
% --- p. 71 ---
\setcounter{footnote}{0}%
\opage{71}immediate instinct of life is too strong for suffering; it will always rise again, --- and even if it succumbs in practice, it will assert itself in thought: thus we see at the close of antiquity the ancient optimism gather its last forces in Plotinus's writing against the pessimism and dualism of the Christian Gnostics.

If we ask more closely what kind of misfortunes can thus penetrate to the marrow of humanity, so that the drive of life is wounded in its innermost part, then \emph{the close of antiquity} presumably offers precisely an illuminating example. When a great and brilliant period of culture is at its consummation, when all the possibilities seem to be used up, when there is, as it were, a pause in world history, then reflection naturally awakens, then the lack of vital force is felt. The old Pan is dead, and one does not suspect that his successor lies in the cradle. No one could have imagined that in a sectarian movement among one of the most insignificant peoples of the Roman Empire, and in the barbarian nations that from time to time stormed the borders of the Empire, lay the germs of a new, fresh and magnificent historical development. The solution of social questions (as well as of many other questions) often lies where one least expects it.

It holds of all development --- as the late Sibbern so often maintained --- that it proceeds sporadically, begins from scattered points, and only when the apparently isolated and accidental series of development converge and unite in the same course does it become possible to foresee the result. Many disappointments arise because one has one-sidedly followed and devoted oneself to a single series of development and staked everything on it, although it can work toward lasting progress only in union with other tendencies. Two of the greatest pessimistic crises of recent times arose in this way, namely those marked by the appearance of Rousseau and of Malthus.

% --- p. 72 ---
\setcounter{footnote}{0}%
\opage{72}Modern culture has its center of gravity in intellectual development. It was in the domain of knowledge that the law of development first dawned clearly upon consciousness. Pascal already uttered the famous words: ``La suite \textit{des hommes, pendant le cours de tant de} siècles, \textit{doit étre considérée comme un} même \textit{homme qui subsiste} toujours \textit{et qui apprend} continuellement'' (``The human race, whose life fills so many centuries, must be regarded as one and the same human being, who always exists and incessantly learns something new.'' \textit{Fragment d'un traité du vide}). Among those who in the eighteenth century worked for progress, Enlightenment therefore naturally became the watchword. Increasing insight was to transform life; in part it had begun to do so. It was to free us from prejudices and antiquated authorities and to bring nature into the service of man. Against those who held that all was well provided people merely learned something, \emph{Rousseau} then set his well-known paradox that art and science were the misfortunes of mankind, and his call to return to the state of nature. Feeling was for him one and all. He has --- as Hettner says --- liberated the human heart and installed the whole and complete human being in its inviolable right, set the forward-pressing language of the feeling heart against the external and rigidly straight-ruled in the one-sided cultivation of the understanding. But it was a sharp judgment he passed on the earlier development, and both he and the minds whose forerunner he became felt painfully the disharmony he had pointed out. Although knowledge and feeling spring from the same root, psychology nevertheless proves that their conditions of development are not the same. Knowledge advances faster than feeling; thought combines more easily with thought than feeling with thought; one thought passes more easily into another thought than one feeling is transformed into another. Feeling thus cannot keep pace. There pass
% --- p. 73 ---
\setcounter{footnote}{0}%
\opage{73}long times, full of strife and pain, especially for those in whom the life of feeling has a peculiar inwardness, before the harmony broken by the swift intellectual development can be restored. It was a great and deep need in human nature, a relation of great psychological and historical significance, that Rousseau pointed to by his pessimistic judgment of modern culture. Here too the natural kinship between pessimism and sympathy showed itself. The optimistic men of the Enlightenment were aristocrats of the mind; their concept of cultivation could be satisfied only by a small circle, beyond which they did not direct their gaze; Rousseau went deeper and therefore also further.\footnote[1]{\emph{Kant}, who in his first writings had paid homage to intellectual optimism, was strongly influenced by Rousseau's writings. He rejected his earlier view as an aristocracy of the mind, and discarded a treatise he had written on optimism. Thereby the foundation was laid for his later ethical view. (See \emph{Dieterich}: Kant und Rousseau. Tübingen 1878. p. 7---13; 97---98).} But in his zeal Rousseau conceived thinking, reflection, as the fruit of a Fall. ``L'état de réflexion est un
état contre la nature, et l'homme qui médite, est un
animal dépravé'' (``To think is contrary to nature, and the human being who ponders is a degenerate creature.''
Discours \textit{sur l'origine de l'inégalité parmi les hommes).}
It is certain that thought leads beyond the state in which the instinct of life rules immediately, in which no possibilities, no comparison, no ideals intervene disturbingly. But Rousseau, this herald of feeling, contradicts himself when, in opposition to reflection, he recommends his so-called state of nature, where ``the imagination paints nothing and the heart demands nothing.'' The point is that although intellectual development for the time being causes disharmony and breaks in upon the instinctive life of feeling, it is
% --- p. 74 ---
\setcounter{footnote}{0}%
\opage{74}nevertheless itself a condition of the higher development of the life of feeling. The height, purity and compass of feeling depend on how many elements of thought it has taken into itself. Hope and fear, love and hatred, admiration and enthusiasm arise only through a fusion of elements of knowledge and elements of feeling. By condemning the life of thought one therefore at the same time deprives a whole series of the highest and noblest feelings of the conditions of their development. Rousseau conceives knowledge only from the side on which it enters into opposition to immediate feeling and, for a time at least, brings strife and unrest in place of the secure repose in instinct. When he demands a return to the state of nature, he demands, really, like the modern pessimists, a return to the unconscious. The presupposition is: Happiness is unconscious; self-consciousness makes us unhappy.

Besides Enlightenment there was another means on which the eighteenth century built its hopes of unlimited progress and happiness for mankind: namely political reforms. If only the external obstacles to universal happiness, which lay in the forms and institutions of society, were removed, all would become well. In its most extravagant form this optimism was voiced by Condorcet and Godwin; but it is prevalent in the literature toward the end of the century (cf. Greg:
\textit{Enigma's of Life.} 4 ed. \textit{London 1873.} p. 54). Then, like a bolt from a clear sky, came Malthus's famous book on the population question, in which he showed that the human race increases faster than the means of subsistence, so that the source of social misfortunes is not to be sought on the surface but in the very foundation of human life. Here too the pessimistic view was the one that went deepest. Since that time social pessimism has several times fought against political optimism, and this disharmony will --- like that kindred
% --- p. 75 ---
\setcounter{footnote}{0}%
\opage{75}psychological disharmony --- require long times for its solution. In both respects it is of the greatest importance that the feeling has been awakened that there is a struggle that must be fought through. --- Social pessimism seems, like the psychological, to require a return to the unconscious. The socialist systems have a more or less reactionary character, place themselves in a more or less polemical relation to individual freedom, which in a way is in society what self-consciousness, intelligence, is in the single human being. And yet it is certain that in the individual freedom for which the struggles of so many ages have paved the way a real good has been attained, which is more than can be said of the social ideals of the future hitherto sketched. To begin by subduing it is just as ill-judged as wanting to save feeling by binding the understanding.

It is not impossible that the strong pessimistic current in \emph{Germany in our day} is for the most part due to this opposition between the political and the social. After the great political end had been attained, which had been the object of so many hopes and exertions, and which had cost so many sufferings, the social disproportions had to emerge all the more sharply. Light makes the darkness blacker. With the attainment of a great end a certain melancholy is also frequently connected. One feels emptiness and asks: what now? That which glittered as the distant end easily loses its luster when it passes into the past. It belongs to the life of feeling to move through oppositions; if barrier and resistance are taken away, the human being feels himself as in empty space. \emph{Stuart Mill} had in his youth a crisis to go through that is of interest in this connection. The view of life into which he had early been led had taught him ``to ground his happiness on a lasting and distant end, toward which
% --- p. 76 ---
\setcounter{footnote}{0}%
\opage{76}progress could constantly be made, without its ever being exhausted by complete attainment.'' But in his enthusiasm for progress and development he was stopped, in a dark hour, by the question: ``Suppose your end in life were fulfilled; suppose all the changes in institutions and opinions that you desire were at this very moment completely realized; --- would you then be inwardly glad and happy?'' And he had to answer this question with No. ``My courage failed me,'' he says; ``the whole foundation on which my life was built collapsed.'' For every view of life there are corresponding trials of faith, in which even that which otherwise stands as the highest and most glorious loses its value and everything becomes empty and desolate. Stuart Mill saw that he had hitherto allowed feeling too small a role; but it was not merely that which led him out beyond the crisis. He was permeated by the conviction that the human being must subordinate himself to an end, a task, in face of which he can forget himself. If one regards oneself as an end, one will feel the emptiness. ``Ask yourself whether you are happy, --- and you will cease to be so. One must not regard happiness, but an end lying beyond it, as the task of life. Let your self-consciousness, your inquiry, your self-observation go out upon this, --- and you will\dots{} breathe in happiness with the air that surrounds you'' (John Stuart Mill's Levned. Copenhagen 1874. p. 136; 145).\footnote[1]{Stuart Mill here approached, as he himself remarks, \emph{Carlyle}'s ``Anti-Self-consciousness Theory.'' For the testimony of this original mind too can be adduced to show that full consciousness --- especially the full consciousness of the attainment of an end --- is the enemy of happiness. ``Self-contemplation,'' he says, ``is always a sign of sickness and egoism. It is unhealthy to consume oneself in empty remorse and fear, and still more unhealthy to wallow in the contemplation of one's own acts. It is of no use to measure the road one has already
\opage{77}covered; what matters is to go forward. If the rational will that dwells in our holy of holies is truly free as a deity, then perfect obedience will be silent. --- --- Unconsciousness and secrecy are marks of everything great; everything one can see completely through is trivial. Rightly did the ancients make silence a god; rightly do more recent poets sing hymns to the night.'' (\emph{Julian Schmidt}: Portraits aus dem 19ten Jahrhundert. 1878. p. 107 f.)} The attainment of an end
% --- p. 77 ---
\setcounter{footnote}{0}%
\opage{77}can make one unhappy by directing reflection toward one's own self, which had hitherto receded before the task. In the endeavor, the advancing striving, happiness must be sought. Therefore Leibnitz could not conceive of bliss otherwise than as constant advancing; otherwise it would bring torpor \textit{(nisi beatitudo in progressu consisteret,
stuperent beati. Epp.} ad Wolf).

German pessimism, however, can scarcely be explained completely by the crisis that the attainment of a great end brings with it. There are certainly many other causes that have contributed to its having come so far that Bismarck last year could say in the German Reichstag: ``The Germans have a strong inclination to discontent; I do not know whether any of us has seen a contented countryman?'' Dr. Waldstein, who has particularly investigated the question of the causes of German pessimism, and by whom an article has been printed in ``Nær og Fjern'' for 1878, holds that the deficiencies and one-sidednesses of both the system of instruction and the sociable and social life have contributed much to this phenomenon of the age. Many of the causes are common to all countries, and could on closer investigation be traced back to the turning points that are marked by Rousseau and Malthus. We are still in the midst of the crises to which these names point. We lack harmony in our view of life, confidence and security with regard to the social foundation of public life.

One must not forget that pessimism's chief
% --- p. 78 ---
\setcounter{footnote}{0}%
\opage{78}spokesman, \emph{Arthur Schopenhauer}, is not a modern author in the narrower sense of the word. His youth falls in the period of Romanticism, the first part of his chief work appeared in 1819, and he was then 31 years old. It has been emphasized from several quarters that pessimistic poetry (Byron, Leopardi, Heine) preceded pessimistic philosophy (Schopenhauer, Hartmann). Pessimistic poetry is closely connected with the direction of feeling issuing from Rousseau. Pessimism's English critic, James Sully, has in a special treatise called attention to the fact that poets at all times, in antiquity too, dwell on the sorrowful side of life. He seeks the cause of this in the fact that it is more important for living beings to express their pain than their pleasure, for which reason the instinct to let pain break out will naturally be the most strongly developed. Now civilized life may well be said to have a tendency to reverse the relation; in good society one must not weep or lament; and moral strength shows itself in one's being able to subdue grief. But the original drive will not be subdued, least of all in the poet, who has a finer, more sensitive constitution than other human beings, and in whom feeling and passion play so great a role. There lies a great consolation and relief in giving suffering a strong and vehement expression. Moreover, the imagination is nourished by the imperfections of the real world; and conversely the world comes precisely to stand as dark and imperfect over against the ideals of fantasy \textit{(Revue philosophique.} Avril 1878).

There can be no sharp boundary drawn between pessimistic poetry and pessimistic philosophy. If by philosophy one understands the critical-scientific investigation of the basic concepts of our thinking, philosophy is neither optimistic nor pessimistic. Its results may at a given time seem to point in the one or the
% --- p. 79 ---
\setcounter{footnote}{0}%
\opage{79}other direction; but this is in itself inessential for it. It is otherwise when one takes the word philosophy in a wider sense. Every view of life that with clear reflection seeks to trace the complicated relations of life back to certain principles may, especially when these principles in turn are sought to be brought into connection with the results of that critical-scientific investigation, be called philosophy. In the systems that the history of philosophy presents to us we usually find both endeavors closely joined, often interwoven with one another. The endeavor of philosophy in our time is to make the critical-scientific investigation the main thing, in order thereby to gain a firm foundation and starting points for secure progress. Schopenhauer occupies in this respect a peculiar position. A great part of his investigations belongs decidedly to scientific philosophy, and to no small degree this has been the ground of his reputation during the reaction against speculative philosophy. He sought to work further on the foundation in the theory of knowledge given by Kant. His philosophizing has an analytic character: from what is given in experience he seeks back to the elements, to the primal phenomena that mark the limits of knowledge. Here he has a firm and secure scientific standpoint over against the speculative philosophy (Fichte, Schelling, Hegel), which wanted to derive everything by a priori thinking from one single absolute principle. But this circumspection does not hold out with him to the end. At a certain point he opens the fence, gives up the exertion and lets the fantasy inspired by feeling rule. ``The thing in itself,'' which Kant had declared to be eternally incomprehensible, we are, according to Schopenhauer, immediately\footnote[1]{Yet Schopenhauer is somewhat wavering here; at times he seems to regard the consciousness of the will as mediate. He says e.g. Parerga und Paralipomena. Berlin 1851. I p. 18: ``Weil doch wir selbst unstreitig real sind, muss \opage{80}aus dem Innern unsers eigenen Wesens die Erkenntniss des Realen irgendwie zu schöpfen sein. In der That nun tritt es hier, \emph{auf eine unmittelbare Weise}, in's Bewusstsein, nämlich als Wille.'' But II p. 79 it says: ``\emph{Wäre dieses Sichbewusstwerden} [des Seienden selbst als Wille] \emph{ein unmittelbares}, so hätten wir eine völlig adäquate Erkenntniss des Dinges an sich. Weil es aber dadurch \emph{vermittelt} ist, dass etc.''} to become conscious,
% --- p. 80 ---
\setcounter{footnote}{0}%
\opage{80}as we, by turning our gaze inward upon ourselves, as it were sense the pulse of existence beat within us, feel the primal force in everything stir as an unruly drive, a yearning and willing that is never fully satisfied. This will, unconscious in itself, is for Schopenhauer the principle of existence; all forces and forms of nature are its expressions, mark steps on its way. Consciousness and thought have arisen merely as a means for the will to attain its end, to satisfy its hunger and thirst for existence. But with consciousness pain also arises. The unconscious will is an ardor that is only provisionally sated, but always announces itself anew. All pleasure is merely negative, is an illusion that arises when the hot desire is quenched for a moment; pain, on the other hand, is positive, the real expression of the primal will when this becomes an object of consciousness. It is therefore above all a matter of abolishing the will to life in order to stop the source of suffering. Various paths stand open to this end: artistic intuition and scientific thought, where man forgets himself in pure contemplation without any desire or willing, --- compassion, in which the individual makes himself wholly one with other individuals and thus abolishes his will, --- resignation and asceticism, which at their height may even lead to man's no longer satisfying the demands of vegetative life. The will thus needs ``quietives,'' not motives; it is to be lulled to sleep like a sick and restless child, not tended and developed.

% --- p. 81 ---
\setcounter{footnote}{0}%
\opage{81}This view of life, which Schopenhauer presents with great energy and poetic gift, scarcely withstands any closer logical examination. In his famous antipodes (whom he never tires of heaping with abuse), in Fichte, Schelling and Hegel, we must precisely admire the power and consistency with which the content of an ideal view of life is formed into a coherent system, with a definite effort to link member to member. But here where their strength lies, precisely there lies Schopenhauer's weakness. The fact is: the deep pessimistic feeling, which in the history of literature so often broke out in passionate outbursts, such as Rousseau in particular and the poets of the first part of the century voiced, without its always being their last word, --- it has become one and all with Schopenhauer, has with him for the first time in the history of thought sought to seize absolute rule. It does so by a coup. For hardly any necessary connection can be shown between the critical-scientific part of his philosophizing and his doctrine of the primal will. With the same general philosophical presuppositions the riddle of the world could just as well have been guessed in another way. It is feeling that has dictated the solution.

But even apart from this, this solution is of great interest. Schopenhauer's view of life stands in decided opposition to the conception of life of the whole of European humanity in its various forms. However deep a feeling of life's pain and sin may have been voiced by founders of religions and philosophical thinkers, the quiet faith nevertheless always underlies it, as a younger poet of ours has expressed in the words that ``above the clouds the sky is always blue.'' That the substance of the world was good, of this one had no doubt; of this even Augustine, with his belief that the majority of men are destined to eternal damnation, agrees with Plato and Leibnitz. And even
% --- p. 82 ---
\setcounter{footnote}{0}%
\opage{82}where one had broken loose from the Platonic and the Christian tradition, one nevertheless believed that the kernel of life, however bitter it might be, could unfold and turn to fruit. But for Schopenhauer rottenness extends all the way into the innermost essence of the world; substance is bad, and all development is an illusion. Optimism is for him a wickedness which the Jews in particular and the modern demagogues have spread. The only thing that reconciled him with the Old Testament was the story of the Fall, and as the real kernel of Christianity he regarded the doctrine of original sin and damnation. On the other hand he took up with great sympathy and enthusiasm the Buddhist doctrine of the pain of existence and of the necessity of returning to Nirvana, the infinite nothing. He feels himself in decided conflict with the European view of life, which in both its religious and its humane forms conceives life as development and progress and places freedom above enjoyment and security. ``The Asiatics,'' says Draper (Den evropæiske Aandsudviklings Historie. Copenhagen 1875. p. 52), ``have weighed freedom and security against each other and deliberately preferred the latter and left it to Europe to sigh after the former. Freedom is prized exclusively in an active life, but the life of Asia is essentially of a passive nature, and it yearns for rest.'' Quietives, not motives, is what Schopenhauer seeks, after all.

It might seem to point to energy in the view of life when Schopenhauer conceives the innermost ground of the world as will. But he is precisely here exceedingly arbitrary and superficial. While in particulars, especially in his more poetic and popular developments, he often betrays fine psychological sense, he lacks in the highest degree precision and thoroughness in determining fundamental psychological concepts. --- He considers it inessential to the will that it is guided by conscious motives; this is only the case with
% --- p. 83 ---
\setcounter{footnote}{0}%
\opage{83}human beings and animals. By contrast it is supposed to be the unconscious will that does everything in nature. The concept of force is to be subordinated to the concept of will; for what force is we understand only by means of the will, of which we are immediately conscious; to understand the unknown we must trace it back to the known. But how then does the will manifest itself in our interior, so that we can become immediately conscious of it? To the will Schopenhauer refers not only all resolutions but also all drives, wishes, passions, and ``all those movements in our interior which one is accustomed to refer under the wide name of feeling.'' He thus does not distinguish between feeling and will. In itself this is indeed a very difficult point in psychology; with a single exception (Alexander Bain) there are no preliminary works for getting clear on this point. But the slackness with which Schopenhauer proceeds is unparalleled.

When one makes feeling and will one, this may come of one's reducing feeling to will, or conversely. Schopenhauer, for whom the will is one with ardor, yearning, drive, conceives it only as passive, thus in reality makes the will into feeling. This is characteristic, since we have seen what an essential role feeling plays every time pessimism appears. Man in Schopenhauer stands predominantly passive toward the external world; he seeks only to cover himself as well as possible and wishes merely to be done with the game as soon as possible. His fighting is at most defensive, never offensive. ``The arrows and slingstones of wanton fate''\footnote[1]{In the famous monologue in which these words occur, Hamlet depicts, as \emph{Bucknill} (\textit{The mad folk of Skakspeare.} p. 86) remarks, his own character: ``so eager to think, so sluggish to act, so keen in his estimate of life's evils, so disinclined to take active steps against them.''} strike him, and
% --- p. 84 ---
\setcounter{footnote}{0}%
\opage{84}he bears it with more or less resignation. The activity of the will must lapse when it is torn loose from its connection with the thought of positive ends that are to be attained. Only through this connection does the will become more than obscure ardor. Fichte, whose ungrateful disciple Schopenhauer was, already made the will into the absolute, but it was precisely as the infinite striving for freedom, the positive working for the realization of the ideal.

But feeling too suffers wrong from that psychological obscurity. Schopenhauer regards pain as the really positive, whereas the feeling of pleasure is always negative, in that it arises when the drive is quenched for a while. How there can be any talk of pleasure at all in Schopenhauer, when the will is the vessel of the Danaids that is never filled, is not easy to see. But in any case he fixes his gaze only on those feelings of pleasure that are immediately tied to self-preservation (e.g. the satisfaction of hunger and thirst). The blind drive clings to life, and when it succeeds, there arises the illusion one calls pleasure, joy or happiness. We rejoice, then, only when we have let ourselves be fooled by the unconscious will. In itself there is no reality in joy or happiness. The greatest goods of life, health, youth and freedom, we prize only when we have lost them.

It is true that in the drive there is frequently a moment of displeasure. But it does not follow that the whole feeling of pleasure in the satisfaction of the drive rests on this displeasure's ceasing. In the drive, so long as it has not decidedly been disappointed, there is at the same time an anticipation, a pleasure in hope, and the feeling of pleasure in the satisfaction of the drive rests in part on this anticipated pleasure's unfolding into full actuality and strength. Even with the drives tied to the maintenance of the individual and of the species there occurs, when they are satisfied, more than when one adds + 3 to --- 3. According to
% --- p. 85 ---
\setcounter{footnote}{0}%
\opage{85}Schopenhauer everything comes out even, and it is the resulting 0 that by an illusion stands for us as positive pleasure. To the normal performance of the organic functions there is tied a happy basic mood, a feeling of lightness and freedom, to which we admittedly seldom turn our attention directly and as a rule notice only when it replaces or is replaced by malaise (just as, according to Helmholtz, we perceive harmony by its contrast to disharmony). But would it be so great a good if we enjoyed our health, youth and freedom with express consciousness? The good in these things consists after all precisely in their putting us in a position for full development and use of our powers. A happy basic mood in the feeling of life is the condition of all our activity; a morbid feeling of life will soon paralyze all our faculties. The feeling of life is the root of our whole mental life; insanity begins as a rule with disturbance and morbidity in it.\footnote[1]{A German physician has sought to show that Schopenhauer suffered from a hereditary insane disposition. His grandmother went insane in her old age. One of her children was an idiot. Schopenhauer's father had a dark and melancholy temperament and suffered from violent fits of anger. Toward the end of his life he showed such an anxiety with regard to his business that he was regarded as insane: his death by drowning in a canal was regarded by many as suicide. --- This anxiety and indefinite fear was also a main feature in Arthur Schopenhauer's temperament. --- (Cf. J. \emph{Sully}: \textit{Pessimism. A history and a criticism. London 1877.} p. 75. 82.)}
Its significance and positive value rests not on its itself pressing forward into the foreground, but on what it bears and conditions.

Schopenhauer overlooks, next, that the feeling of pleasure need not have pain as its background, even though it stands out most sharply in this relation of contrast. It can
% --- p. 86 ---
\setcounter{footnote}{0}%
\opage{86}have a weaker feeling of pleasure as its background. That would surely be the happiest existence, where preceding joys could serve the following as foil. Barriers to our activity arise not only by an external power's intervening in our existence and causing positive pain, but also by there being still something that has not been attained, though it can be attained. These barriers can claim us so considerably that, without actual preceding displeasure, a positive pleasure arises in the carrying through of the work. In the very striving, so long as it advances, --- so long as the resistance does not exceed a certain limit, --- the moment of pleasure is predominant. But this is an ideal that only in part falls to our lot. Our conditions of life are such that pain incessantly intervenes with its dark shadows. And precisely in more deeply developed individuals this will be the case, since they bear not only their own sufferings, but in sympathy also those of others. Yet it is certain that precisely the noblest enjoyments are raised above the unrest and pain of the drive of self-preservation. Schopenhauer must concede this regarding the joy that art and thought cause; but he thereby comes into contradiction with himself: for in pure Nirvana --- and only there is ``the will to life'' completely abolished --- there can be neither intuition of art nor thought.

But Schopenhauer basically cannot explain the difference between ideal and physical enjoyments. As he recognizes no historical development in the species, so neither does he recognize any psychological development in the individual. That a life of thought develops in man is, in his opinion, only to bring it about that man is shaken far more strongly by the passions than the animal; but it is not to exert any essential influence upon the drives and feelings themselves: ``\emph{am Ende und im Realen} handelt es sich nur um die \emph{selben Dinge}, die auch das Thier erlangt, und zwar mit unvergleichlich geringerem Aufwande von Affekten
% --- p. 87 ---
\setcounter{footnote}{0}%
\opage{87}und Quaalen'' (Parerga und Paralipomena. Berlin 1851. II p. 250 f). This is consistent, when the will as blind insatiable drive stirs in everything. With increasing consciousness everything must only grow worse. But it conflicts entirely with psychological observation, which precisely (as already Plato and the Stoics have shown) demonstrates the purifying, developing and harmonizing influence of thought on the life of feeling. Under the inward reciprocal relation that psychology shows between knowledge and feeling, the former borrows its strength from the latter, but gives it in return content, scope and clarity. As we have seen above, it costs long times before the development of knowledge fully benefits feeling; but that it can happen, and that the natural development leads in that direction, takes great blindness to deny. It is possible that the brutal instincts of self-preservation (Schiller's Hunger und Liebe) still essentially dominate life; but it is in any case a significant progress that they are now compelled to demonstrate their justification by being tied to a more or less ideal and universal end (in the family, the state, etc.). The ideal, which from the first is means, becomes gradually end.

The a priori assumption of blind drive and ardor as the kernel of life, an assumption which the pessimistic feeling itself dictated, makes it impossible for Schopenhauer to be an impartial observer of reality. One would think that experience was the only way by which one could learn whether there is happiness in the world. Out of one's own interior one cannot fetch the decision of this question; for because I am unhappy, it is not said that all others are. Schopenhauer is soon done with experience. ``If one,'' he says, ``wants briefly to test the assertion that enjoyment outweighs suffering in this world, or that they at least balance each other, then
% --- p. 88 ---
\setcounter{footnote}{0}%
\opage{88}compare the feeling of the animal that devours another with that other animal's feeling.'' There is of course greater pain in being eaten than there is joy in eating. But the sum is not so simple. Human beings have after all at least secured themselves tolerably against being eaten; their life need not be a constant fear. But the testimony of experience Schopenhauer cannot attach great weight to. Even if one would assure him that one was happy, he would deny that it was true. He knows, after all, what lies behind it, --- that it is ``the will to life'' we blindly follow. But it is a presumption to call all pleasure an illusion. Whether pleasure or pain is really felt, only the individual can decide, and he does not let himself be disputed out of his feeling. Therefore, e.g., reasoning helps so little against hypochondria, ``reason'' so little against passion. Even the so-called illusory, untrue or imagined joy is a real joy. What one means by these expressions is actually that the cause of the feeling is an imagination or has no subsistence. Only by criticizing the cause and content of the feeling can one get at it itself.

The older pessimists still conceived the unconscious as happy. The happiness they did not find in the present they put back into a golden age, a paradise or a childlike state of nature. But modern pessimism, which fancies it has discovered the innermost of existence, conceives the unconscious life as pain and unrest. It thereby entangles itself in one of its many contradictions; for on the other hand precisely a return to the unconscious is to be the remedy for all pain.

Schopenhauer's successor, Eduard v. Hartmann, has in several essential respects corrected his doctrine. On several points (e.g. in the question of the
% --- p. 89 ---
\setcounter{footnote}{0}%
\opage{89}immediate conception of the will) he admittedly shows greater clarity and prudence. But he is a pessimist only in theory, or --- as one of Schopenhauer's more faithful adherents, \emph{Julius Bahnsen}, has said --- one does not know whether to call him a pessimistic optimist or an optimistic pessimist (Zur Philosophie der Geschichte. Berlin 1872. p. 3). Hartmann namely acknowledges the reality of historical development; it is for him more than frost flowers on the pane or figures in the clouds, to use Schopenhauer's images. But as the final end of development he then posits insight into the wretchedness of life and a resolution arising from it to stop it. What with Schopenhauer is each individual's task, falls to the whole species --- and is thereby shot out into an indeterminate future. In the meantime the optimistic illusions then play their game; the end comes only when they have all run out the line. Hartmann stands to Schopenhauer as the modern Christians to the believers of the apostolic age: these expected Judgment Day in the near future; for those it lies out beyond the horizon, so that one can for the time being arrange things rather cozily. It is therefore characteristic that Hartmann invites his reader to visit him to see how happily one lives in the pessimist's home.

Hartmann has sought to make good what Schopenhauer neglected, to go through the individual sources of enjoyment and examine their value. But he too, by virtue of his exact knowledge of ``the unconscious,'' pretends to know better about others' happiness or unhappiness than their own feeling. Moreover --- as \emph{James Sully} in his book cited above has shown --- his investigation is highly superficial. It is especially remarkable that he does not dwell on the feeling of pleasure that springs from activity and work, but mainly keeps to the passive enjoyments. Compelled or of free will, in earnest or as play, we spend after all the greatest
% --- p. 90 ---
\setcounter{footnote}{0}%
\opage{90}part of our life in movement and activity. If self-reflection makes unhappy, it is precisely because one stops on the way, draws thought away from the end, the task, in which it can find rest. What is true in the doctrine of the unconscious as a universal remedy is not that we should go back, but precisely that we should go forward. Self-forgetfulness is precisely to serve to develop mental life further. The unconscious is not a nightmare that weighs on us, but a source from which life ever draws new strength.

Schopenhauer and Hartmann share the illusion of speculative philosophy that thought can posit an absolute principle. Life and experience show us only struggling forces, a play of oppositions, and it is arbitrary to make any single one of these into everything. On the other hand, these forces and oppositions have a highly different value for human life and its higher development. Even the most valuable exists here in the world only as a relative, struggling power. He who would win it gets a share in it only by having a share in its struggle to subsist and develop. Neither optimism nor pessimism corresponds to reality when they are to be carried through consistently. Meliorism, an expression proposed by George Elliot (Sully: Pessimism. p. 399), will be the right word when by it one understands a belief in the possibility of advance and development. And when we hold fast the experience which psychology and history no less than external nature impress upon us, that development tends to proceed from scattered points, then there presents itself both the task and the hope that the elements, in their separation contending, will by the law of development be formed and united into wholeness and harmony.
\signature{\textbf{\sffamily Harald Høffding.}}

\end{document}
